\chapter{1977年重庆卷}

\section{解答题}

\begin{problem}
解方程：\(9\left(3x+2\right)^2-4=0\).
\end{problem}

\begin{solution}
移项得
\[
9\left(3x+2\right)^2=4
\]
将左边看作平方差，得
\[
\left[3\left(3x+2\right)+2\right]\left[3\left(3x+2\right)-2\right]=0
\]
即
\[
\left(9x+8\right)\left(9x+4\right)=0
\]
所以 \(9x+8=0\) 或 \(9x+4=0\)，从而
\[
x=-\frac{8}{9}\quad\text{或}\quad x=-\frac{4}{9}
\]
把这两个值分别代回原方程，均使等式成立，故原方程的解为
\(x_1=-\frac{8}{9}\)，\(x_2=-\frac{4}{9}\).

也可由 \(\left(3x+2\right)^2=\frac49\) 得
\(3x+2=\pm\frac23\)，解得同样的两个值.
\end{solution}

\begin{problem}
已知 \(\lg5=0.6990\)，问 \(8^{34}\times25^{11}\) 是多少位的整数.
\end{problem}

\begin{solution}
设 \(N=8^{34}\times25^{11}\). 由 \(\lg10=1\) 和已知条件，
\[
\lg2=\lg10-\lg5=1-0.6990=0.3010
\]
于是
\[
\begin{aligned}
\lg N
&=34\lg8+11\lg25\\
&=34\times3\lg2+11\times2\lg5\\
&=102\times0.3010+22\times0.6990\\
&=30.702+15.378=46.0800.
\end{aligned}
\]
因此 \(10^{46}<N<10^{47}\). 正整数的常用对数在 \((46,47)\) 内时，该正整数恰有 \(47\) 位，所以 \(8^{34}\times25^{11}\) 是 \(47\) 位的整数.
\end{solution}

\begin{problem}
把 \(\sin\left(30^\circ+x\right)+\cos\left(x+60^\circ\right)-\sin^2x-\cos^2x\) 化成积的形式.
\end{problem}

\begin{solution}
利用 \(\cos\theta=\sin\left(90^\circ-\theta\right)\) 和平方关系，原式为
\[
\sin\left(30^\circ+x\right)+\sin\left(30^\circ-x\right)-1
\]
用两角和的正弦之和公式，得
\[
\begin{aligned}
&\sin\left(30^\circ+x\right)+\sin\left(30^\circ-x\right)-1\\
&=2\sin30^\circ\cos x-1\\
&=\cos x-1\\
&=-\left(1-\cos x\right)\\
&=-2\sin^2\frac{x}{2}.
\end{aligned}
\]
所以原式化成积的形式为 \(-2\sin^2\frac{x}{2}\).
\end{solution}

\begin{problem}
试证明：\(\left[\left(2m\right)^0-\left(m+x\right)^{-2}\right]\div\left[\left(2x\right)^0-\left(m+x\right)^{-1}\right]^2\cdot\left[2m^0-\frac{1-\left(m-x\right)^2}{2mx}\right]=\frac{\left(m+x+1\right)^2}{2mx}\).
\end{problem}

\begin{solution}
原式有意义并且可以进行下列约分时，需要
\(m\ne0\)，\(x\ne0\)，\(m+x\ne0\)，\(m+x\ne1\). 在这些条件下，先把幂改写为分式：
\(\left(2m\right)^0=1\)，\(\left(2x\right)^0=1\)，\(\left(m+x\right)^{-2}=\frac{1}{\left(m+x\right)^2}\)，\(\left(m+x\right)^{-1}=\frac{1}{m+x}\).
因此第一因子为
\[
1-\frac{1}{\left(m+x\right)^2}
=\frac{\left(m+x-1\right)\left(m+x+1\right)}{\left(m+x\right)^2}
\]
被除数中的平方因子为
\[
\left(1-\frac{1}{m+x}\right)^2
=\frac{\left(m+x-1\right)^2}{\left(m+x\right)^2}
\]
第三因子化为
\[
\begin{aligned}
2-\frac{1-\left(m-x\right)^2}{2mx}
&=\frac{4mx-1+\left(m-x\right)^2}{2mx}\\
&=\frac{m^2+x^2+2mx-1}{2mx}\\
&=\frac{\left(m+x-1\right)\left(m+x+1\right)}{2mx}.
\end{aligned}
\]
代入原式并约去非零因子，得
\[
\begin{aligned}
&\frac{\left(m+x-1\right)\left(m+x+1\right)}{\left(m+x\right)^2}
\cdot\frac{\left(m+x\right)^2}{\left(m+x-1\right)^2}
\cdot\frac{\left(m+x-1\right)\left(m+x+1\right)}{2mx}\\
&=\frac{\left(m+x+1\right)^2}{2mx}.
\end{aligned}
\]
这正是题中右端，故等式成立.
\end{solution}

\begin{problem}
将半径为 \(R\)，中心角为 \(120^\circ\) 的扇形铁片，卷成一圆锥容器，问圆锥容器的容积是多少（接缝及卷边不计）.
\end{problem}

\begin{solution}
设卷成的圆锥底面半径为 \(r\)，高为 \(h\). 扇形的两条半径卷合后成为圆锥的母线，所以母线长仍为 \(R\)，扇形弧长成为圆锥底面周长.

扇形弧长为
\[
\frac{120^\circ}{360^\circ}\cdot2\pi R=\frac{2\pi R}{3}
\]
所以
\(2\pi r=\frac{2\pi R}{3}\)，\(r=\frac{R}{3}\).
由圆锥的母线、高和底面半径构成直角三角形，
\[
h=\sqrt{R^2-r^2}
=\sqrt{R^2-\left(\frac{R}{3}\right)^2}
=\frac{2\sqrt2}{3}R
\]
于是圆锥容积为
\[
V=\frac13\pi r^2h
=\frac13\pi\left(\frac{R}{3}\right)^2\cdot\frac{2\sqrt2}{3}R
=\frac{2\sqrt2}{81}\pi R^3
\]
故圆锥容器的容积为 \(\frac{2\sqrt2}{81}\pi R^3\).

\solutionfigure{\bitmapfigure[max width=3.2cm]{img/1977/chongqing/q2_solution_fig1.png}}
\end{solution}

\begin{problem}
如图，\(AD\)，\(BD\) 分别为 \(\triangle ABC\) 的角平分线，延长 \(AD\) 交 \(\triangle ABC\) 的外接圆于 \(E\)，连接 \(BE\)，求证：\(BE=DE\).

\FigureLayoutDeclare{img/1977/chongqing/q3_fig1.png}{10.0cm}{0bp 0bp 0bp 0bp}
\bitmapfigure[max width=6.0cm]{img/1977/chongqing/q3_fig1.png}
\end{problem}

\begin{solution}
延长 \(BD\) 交外接圆于 \(F\). 由于两弦 \(AE\) 和 \(BF\) 在圆内相交，圆内两弦所成角的度数等于所对两弧度数和的一半，因此
\[
m\left(\angle BDE\right)=\frac12\left[m\left(\widehat{BE}\right)+m\left(\widehat{AF}\right)\right]
\]
又因为 \(D\) 在 \(BF\) 上，\(\angle DBE=\angle FBE\) 是圆周角，它所对的弧为不含点 \(B\) 的弧 \(FE\)，所以
\[
m\left(\angle DBE\right)=\frac12m\left(\widehat{FE}\right)
\]
图中弧 \(FE\) 经过 \(C\)，故
\[
m\left(\angle DBE\right)=\frac12\left[m\left(\widehat{EC}\right)+m\left(\widehat{CF}\right)\right]
\]
因为 \(AD\) 平分 \(\angle BAC\)，且 \(A\)，\(D\)，\(E\) 共线，所以
\(\angle BAE=\angle EAC\). 这两个圆周角分别所对弧 \(BE\) 和 \(EC\)，从而
\[
m\left(\widehat{BE}\right)=m\left(\widehat{EC}\right)
\]
同理，\(BD\) 平分 \(\angle ABC\)，且 \(B\)，\(D\)，\(F\) 共线，所以
\(\angle ABF=\angle FBC\)，从而
\[
m\left(\widehat{AF}\right)=m\left(\widehat{FC}\right)
\]
代入前面的两个角度关系，得到
\[
m\left(\angle BDE\right)=m\left(\angle DBE\right)
\]
所以 \(\triangle BDE\) 的两个底角相等，故 \(BE=DE\).
\end{solution}

\begin{problem}
顺次连接 \(A(-2,4)\)，\(B(2,2)\)，\(C(4,0)\) 得 \(\triangle ABC\)，写出 \(BC\) 边上的高所在的直线方程，并计算 \(\triangle ABC\) 的面积.
\end{problem}

\begin{solution}
直线 \(BC\) 的斜率为
\[
\frac{0-2}{4-2}=-1
\]
所以直线 \(BC\) 的方程为
\(y-2=-(x-2)\)，\(x+y-4=0\).
\(BC\) 边上的高垂直于 \(BC\)，其斜率为 \(1\)，又经过点 \(A(-2,4)\)，故高所在直线为
\(y-4=x+2\)，\(y=x+6\).

点 \(A\) 到直线 \(BC\) 的距离为高的长度：
\[
h=\frac{|(-2)+4-4|}{\sqrt{1^2+1^2}}=\sqrt2
\]
而
\[
|BC|=\sqrt{(4-2)^2+(0-2)^2}=2\sqrt2
\]
因此
\[
S_{\triangle ABC}=\frac12h|BC|=\frac12\cdot\sqrt2\cdot2\sqrt2=2
\]
所以高所在直线方程为 \(y=x+6\)，三角形面积为 \(2\).
\end{solution}

\begin{problem}
如图，在梯形 \(ABCD\) 中，\(a:b:c:d=1:3:1:2\). 求在梯形 \(ABCD\) 中，\(\cos\alpha\)，\(\sin\alpha\) 和 \(\sin\beta\) 的值.

\bitmapfigure[max width=4.8cm]{img/1977/chongqing/q5_fig1.png}
\end{problem}

\begin{solution}
由图记 \(AB=a\)，\(DC=b\)，\(AD=c\)，\(BC=d\)，且 \(AB\parallel DC\). 从 \(A\)、\(B\) 分别向 \(DC\) 作垂线，垂足为 \(E\)、\(F\). 因为两条垂线之间的四边形 \(ABFE\) 是矩形，所以
\(EF=AB=a\)，\(AE=BF\).
设 \(DE=x\)，则
\[
FC=DC-DE-EF=b-x-a=b-a-x
\]
由比例条件可取 \(a\) 为基本长度，于是
\(c=a\)，\(b=3a\)，\(d=2a\).
在直角三角形 \(ADE\) 和 \(BFC\) 中，勾股定理分别给出
\(AE^2=a^2-x^2\)，\(BF^2=(2a)^2-(3a-a-x)^2\).
利用 \(AE=BF\)，得
\[
a^2-x^2=4a^2-(2a-x)^2=4ax-x^2
\]
所以 \(4ax=a^2\)，即
\(DE=x=\frac a4\)，\(qquad FC=2a-x=\frac{7a}{4}\).
在直角三角形 \(BFC\) 中，\(\alpha\) 位于点 \(C\)，因此
\[
\cos\alpha=\frac{FC}{BC}=\frac{\frac{7a}{4}}{2a}=\frac78
\]
\(BF=\sqrt{(2a)^2-\left(\frac{7a}{4}\right)^2}=\frac{\sqrt{15}}4a\)，\(\sin\alpha=\frac{BF}{BC}=\frac{\sqrt{15}}8\).
在直角三角形 \(ADE\) 中，\(\beta\) 位于点 \(D\)，所以
\[
\sin\beta=\frac{AE}{AD}=\frac{BF}{c}
=\frac{\frac{\sqrt{15}}4a}{a}=\frac{\sqrt{15}}4
\]
故 \(\cos\alpha=\frac78\)，\(\sin\alpha=\frac{\sqrt{15}}8\)，\(\sin\beta=\frac{\sqrt{15}}4\).
\end{solution}

\begin{problem}
分别在三个平面直角坐标上作出下列方程的图形.

\begin{enumerate}
\item \(\left(x+2\right)\left(y-1\right)=0\)；
\item \(\left(x+2\right)^2+\left(y-1\right)^2=0\)；
\item \(\left(x+2\right)^2-\left(y-1\right)^2=0\).
\end{enumerate}
\end{problem}

\begin{solution}
\begin{enumerate}
\item 乘积为零，等价于至少有一个因式为零，因此
\[
x+2=0\quad\text{或}\quad y-1=0
\]
所以图形是两条直线 \(x=-2\) 和 \(y=1\).

\solutionfigure{\bitmapfigure[max width=4.0cm]{img/1977/chongqing/q6_fig1.png}}

\item 两个平方数都不小于零，平方数之和等于零，当且仅当两个平方数同时为零. 因此
\[
\begin{cases}
x+2=0,\\
y-1=0,
\end{cases}
\qquad\text{即}\qquad
\begin{cases}
x=-2,\\
y=1.
\end{cases}
\]
所以图形只有一个点 \((-2,1)\).

\solutionfigure{\bitmapfigure[max width=4.0cm]{img/1977/chongqing/q6_fig2.png}}

\item 将平方差分解，得
\[
\begin{aligned}
\left(x+2\right)^2-\left(y-1\right)^2
&=\left(x+2-y+1\right)\left(x+2+y-1\right)\\
&=\left(x-y+3\right)\left(x+y+1\right).
\end{aligned}
\]
因此原方程等价于
\[
x-y+3=0\quad\text{或}\quad x+y+1=0
\]
所以图形是两条直线 \(x-y+3=0\) 和 \(x+y+1=0\).

\solutionfigure{\bitmapfigure[max width=4.0cm]{img/1977/chongqing/q6_fig3.png}}
\end{enumerate}
\end{solution}

\section{参考题}

\begin{problem}
求函数 \(f(x)=\frac{1}{4}x^4-x^3-2x^2+3\) 的极值.
\end{problem}

\begin{solution}
函数在 \(\R\) 上可导，先求其可能的极值点：
\[
f'(x)=x^3-3x^2-4x=x(x-4)(x+1)
\]
令 \(f'(x)=0\)，得临界点 \(x=-1\)，\(0\)，\(4\). 再求二阶导数：
\[
f''(x)=3x^2-6x-4
\]
于是
\(f''(-1)=5>0\)，\(f''(0)=-4<0\)，\(f''(4)=20>0\).
由二阶导数判别法，\(x=-1\) 和 \(x=4\) 处为极小值点，\(x=0\) 处为极大值点.

分别计算函数值：
\[
f(-1)=\frac14+1-2+3=\frac94=2\frac14
\]
\[
f(0)=3
\]
\[
f(4)=\frac14\times4^4-4^3-2\times4^2+3=64-64-32+3=-29
\]
所以当 \(x=0\) 时有极大值 \(3\)；当 \(x=-1\) 时有极小值 \(2\frac14\)；当 \(x=4\) 时有极小值 \(-29\).
\end{solution}

\begin{problem}
计算 \(\int\left(\frac{1}{\sqrt{1-x^2}}-\frac{x^2}{\sqrt{4-x^2}}\right)\mathrm{d}x\).
\end{problem}

\begin{solution}
实被积函数在 \((-1,1)\) 上有意义. 第一项直接积分得
\[
\int\frac{1}{\sqrt{1-x^2}}\mathrm{d}x=\arcsin x+C_1
\]
对第二项令 \(x=2\sin t\). 在本题的实数定义域内可以取 \(t=\arcsin\frac{x}{2}\)，且 \(\cos t>0\)，于是
\(\mathrm{d}x=2\cos t\,\mathrm{d}t\)，\(\sqrt{4-x^2}=2\cos t\).
因此
\[
\begin{aligned}
\int\frac{x^2}{\sqrt{4-x^2}}\mathrm{d}x
&=\int\frac{4\sin^2t}{2\cos t}\cdot2\cos t\,\mathrm{d}t\\
&=\int4\sin^2t\,\mathrm{d}t\\
&=\int2\left(1-\cos2t\right)\,\mathrm{d}t\\
&=2t-\sin2t+C_2.
\end{aligned}
\]
代回 \(t=\arcsin\frac{x}{2}\)，并利用
\[
\sin\left(2\arcsin\frac{x}{2}\right)
=2\cdot\frac{x}{2}\cdot\sqrt{1-\frac{x^2}{4}}
=\frac{x}{2}\sqrt{4-x^2}
\]
得到
\[
\int\frac{x^2}{\sqrt{4-x^2}}\mathrm{d}x
=2\arcsin\frac{x}{2}-\frac{x}{2}\sqrt{4-x^2}+C_2
\]
所以原积分为
\[
\begin{aligned}
&\int\left(\frac{1}{\sqrt{1-x^2}}-\frac{x^2}{\sqrt{4-x^2}}\right)\mathrm{d}x\\
&=\arcsin x-\left(2\arcsin\frac{x}{2}-\frac{x}{2}\sqrt{4-x^2}\right)+C\\
&=\arcsin x-2\arcsin\frac{x}{2}+\frac{x}{2}\sqrt{4-x^2}+C.
\end{aligned}
\]
\end{solution}

\begin{problem}
计算 \(\iint_D\left(1+y^2\right)\mathrm{d}x\mathrm{d}y\)，区域 \(D\) 是 \(y=x^2\)，\(x=y^2\) 所围成的区域.
\end{problem}

\begin{solution}
联立两条边界曲线. 由 \(y=x^2\) 和 \(x=y^2\) 得 \(x=x^4\)，所以交点的横坐标为 \(x=0\) 或 \(x=1\)，相应交点为 \((0,0)\) 和 \((1,1)\). 在 \(0\leqslant x\leqslant1\) 内，\(y=x^2\) 是下边界，\(y=\sqrt{x}\) 是上边界，因此由二重积分的累次积分和 Fubini 定理，
\[
\iint_D\left(1+y^2\right)\mathrm{d}x\mathrm{d}y
=\int_0^1\left(\int_{x^2}^{\sqrt{x}}\left(1+y^2\right)\mathrm{d}y\right)\mathrm{d}x
\]
先对 \(y\) 积分：
\[
\begin{aligned}
\int_{x^2}^{\sqrt{x}}\left(1+y^2\right)\mathrm{d}y
&=\left[y+\frac{y^3}{3}\right]_{x^2}^{\sqrt{x}}\\
&=\sqrt{x}+\frac{x\sqrt{x}}{3}-x^2-\frac{x^6}{3}.
\end{aligned}
\]
再对 \(x\) 积分，得
\[
\begin{aligned}
\iint_D\left(1+y^2\right)\mathrm{d}x\mathrm{d}y
&=\left[-\frac{x^7}{21}-\frac{x^3}{3}+\frac{2x^2\sqrt{x}}{15}+\frac{2x\sqrt{x}}{3}\right]_0^1\\
&=-\frac1{21}-\frac13+\frac2{15}+\frac23=\frac{44}{105}.
\end{aligned}
\]
\end{solution}

\begin{problem}
求下列方程的通解.

\begin{enumerate}
\item \(x\mathrm{d}y-2y\mathrm{d}x=3x^2\mathrm{d}x\)；
\item \(y''-2y'+y=\cos x\).
\end{enumerate}
\end{problem}

\begin{solution}
\begin{enumerate}
\item 对于第一式，先在 \(x\ne0\) 的区间上把方程两边同除以 \(x^3\)，得
\[
\frac{1}{x^2}\mathrm{d}y-\frac{2y}{x^3}\mathrm{d}x=\frac{3}{x}\mathrm{d}x
\]
注意到
\(\mathrm{d}\left(\frac{y}{x^2}\right)=\frac{1}{x^2}\mathrm{d}y-\frac{2y}{x^3}\mathrm{d}x\)，\(\mathrm{d}\left(3\ln|x|\right)=\frac{3}{x}\mathrm{d}x\).
所以
\[
\mathrm{d}\left(\frac{y}{x^2}\right)=\mathrm{d}\left(3\ln|x|\right)
\]
积分得
\[
\frac{y}{x^2}=3\ln|x|+C
\]
从而通解为
\(y=3x^2\ln|x|+Cx^2\)，\(x\ne0\).

\item 先求对应齐次方程
\(y''-2y'+y=0\) 的通解. 其特征方程为
\[
\lambda^2-2\lambda+1=(\lambda-1)^2=0
\]
所以有二重特征根 \(\lambda=1\)，齐次通解为
\[
y_h=\left(C_1+C_2x\right)\e^x
\]
为求非齐次方程的一个特解，设
\[
y_p=A\cos x+B\sin x
\]
则
\(y_p'=-A\sin x+B\cos x\)，\(y_p''=-A\cos x-B\sin x\).
代入左端，得到
\[
y_p''-2y_p'+y_p=-2B\cos x+2A\sin x
\]
要使其等于 \(\cos x\)，必须有
\(-2B=1\)，\(2A=0\)
即 \(A=0\)，\(B=-\frac12\). 因而
\(y_p=-\frac12\sin x\)，原方程的通解为
\[
y=\left(C_1+C_2x\right)\e^x-\frac12\sin x
\]
\end{enumerate}
\end{solution}
